\mychapter{2}{Present state of CDR}
This section presents an overview of the present state of carbon removal, based on the recent developments of  carbon markets, literature on the global governance mechanisms of carbon removal, and the four negative emission technologies that were selected based on their maturity and prevalence, both in terms of research literature and actual deployment.

\section{Relevance of CDR and Carbon Markets}
 The \textcite{unep_emissions_2017} expects that, in addition to stringent emissions reductions, about 10 Gt of annual CDR will be required by 2050, rising to 20 Gt annually by 2100. \textcite{breyer_proposing_2023} estimate the overall CDR required to achieve the 1.5 C target to be about 500 Gt over the course of the 21st century.\vspace{2mm}

In comparison, according to \textcite{cdrfyi_2025_2025}, only about 15 Mt of CDR were contracted, but not yet delivered, in Q2 of 2025. Of those, more than 90\% were contracted by a single corporate buyer and are largely based on BECCS projects. This aligns with the findings by \textcite{roston_market_2025}, who argues that the voluntary CDR markets, despite being the primary driver of early-stage CDR projects, suffer from a lack of standardized products, face liquidity issues, and, through dependence on few, but large transactions, prevent efficient price discovery. It is therefore imperative to move beyond reliance on present VCMs to understand the optimal composition of future CDR portfolios, and instead analyze the available approaches based on their (projected) availability, cost and ultimately their ability to achieve an increase of annual CDR by at least two orders of magnitude until 2050.

\section{Defining CDR}
There is a wide range of terminology and definitions in the carbon removal industry. The most common terms include carbon dioxide removal (CDR) and synonyms such as greenhouse gas removal (GGR) or simply carbon removal. GGR technically includes the removal of non-CO$_2$ greenhouse gases (GHGs), however, none of these can presently be removed at scale \parencite{net_zero_climate_greenhouse_nodate}. The different approaches used to conduct CDR are typically referred to as negative emission technologies (NET's), irrespective of whether it is a natural, engineered or hybrid approach. Finally, the use of NET's to conduct CDR yields what is typically referred to as a "carbon credit" - the claim on one tonne of CO$_2$ removed from the atmosphere.\vspace{2mm}

\textcite{smith_state_2024} defines CDR by the three principles of (1) removing carbon from the atmosphere, (2) using durable storage, and (3) as the result of active human intervention. These principles align with the \textcite{ipcc_global_2019} report on global warming, which defines CDR as "anthropogenic activities removing CO$_2$ from the atmosphere and durably storing it in geological, terrestrial or ocean reservoirs, or in products." This definition is crucial, since it clarifies that point capture (CCUS), e.g. at fossil power plants, and avoidance projects aimed on the prevention of GHG emissions, do not qualify as CDR. It is also intuitive, considering that the prevention of emissions results in a net-zero effect, and using such credits to offset additional carbon emissions would result in double-counting \parencite{reichelstein_accounting_2024}.\vspace{2mm}

The cost of carbon removal is typically referred to as the levelized cost of carbon captured (LCOC). The baseline LCOC is calculated from the net operating costs and capital expenditures divided by the net carbon removed. From the base LCOC, a number of further terms are derived. The LCOC$^+$ includes the leakage from the carbon removal process, e.g. emissions from energy use, which is particularily relevant for direct air capture. The LCOCR additionally includes the transport and storage costs, while the LCOCR$^+$ also considers the leakage from production, transport and storage. Finally the LCOCR$^{+*}$ deducts any co-revenues generated, for example from the energy produced by BECCS.

\section{CDR Governance}
The governance of carbon removal, and global carbon emissions in general, is still lacking, making the comparison of different carbon credit offerings difficult, and the existing carbon markets highly in-transparent \parencite{roston_market_2025}.\vspace{2mm}

A useful way of understanding this issue is to draw a comparison to money: Economists generally follow the philosophy of "a dollar is a dollar is a dollar" - the fungibility assumption declares that each dollar (or Euro, or any other currency) is a perfect substitute for any other dollar \parencite{zelizer_dollar_2017}. And fundamentally, this assumptions should also apply to CO$_2$: Each tonne of carbon dioxide emitted to the atmopshere is as good (or rather, bad) as any other tonne anywhere else in the world. The global warming effect is always the same \parencite{ritchie_co_2023}.\vspace{2mm}

However, in the carbon removal market, as the vast price differences for carbon credits show, the fungibility assumption obviously fails \parencite{roston_market_2025}. Otherwise, why would someone pay 500 US\$ for a carbon credit, if another is available for a mere 4 US\$? \parencite{reichelstein_accounting_2024} The root cause of these discrepancies lies in the fact that while a tonne of carbon emitted is always worth the same (negative) value, a tonne removed is decidedly not. Instead, the quality, and thus economic value, of a carbon credit is defined by its (1) verifiability, (2) additionality and (3) permanence. Other considerations also include accounting, monitoring, or avoidance of negative side-effects, but these can largely be summarized under the former three categories \parencite{axelsson_revised_2024}.\vspace{2mm}

The first two categories, verifiability and additionality, are both questions of quantification: How can we ensure, that any given carbon credit is in fact backed by real carbon removed from the atmosphere, and that this carbon removed amounts to one net tonne? On the other hand, permanence is a question of storage: How can we ensure that this tonne of CO$_2$ is stored in such a manner that it will stay out of the atmosphere for the promised period of time? \parencite{loffler_2025_2025}\vspace{2mm}

There are a number of different frameworks that try to solve these governance problems, such as Gold Standard, Verra, or the GCC regulatory framework, among many others. However, an analysis of green-washing scandals and VCM failures by \textcite{sasaki_addressing_2025} showed that even those frameworks have not been able to consistently ensure the quality of  carbon credits certified through them. More recently, governments have also started to adopt their own CDR verification frameworks. The European Union's "Carbon Removals and Carbon Farming" regulation aims to establish an EU-wide taxonomy of NET's considered permanent, and standards to ensure quality, monitoring, and reporting, while, importantly, facilitating investments in novel carbon removal technologies \parencite{european_parliament_regulation_2024}.

\section{CDR Approaches and Maturity}
The literature broadly categorizes CDR approaches into nature-based solutions, engineered ("novel"), and hybrid solutions. Each category presents specific trade-offs between the cost, scalability, and quality of the removals.

\begin{table}[]
\small
\renewcommand{\arraystretch}{0.9}
\begin{tabular}{lllll}
Technology     & TRL   & Primary Cost Drivers                                                                & Benefits                     & Drawbacks                                                                            \vspace{1.5mm}\\ \hline & \\[-1.5ex]
\textbf{A/R}   & 8 - 9 & \begin{tabular}[c]{@{}l@{}}Land acquisition, planting,\\ monitoring\end{tabular}    & Proven, cheap                & \begin{tabular}[c]{@{}l@{}}Limited permanence,\\ difficult to quantify\end{tabular}  \vspace{1.5mm}\\
\textbf{ERW}   & 4 - 5 & \begin{tabular}[c]{@{}l@{}}Processing energy, transport,\\ spreading\end{tabular}   & Gigatonne-scale availability & \begin{tabular}[c]{@{}l@{}}Mining logistics,\\ socio-political barriers\end{tabular} \vspace{1.5mm}\\
\textbf{BiCRS} & 7 - 9 & \begin{tabular}[c]{@{}l@{}}Feedstock supply, BECCS\\ CCS units\end{tabular}         & Energy, soil improvement     & Biomass Supply, (Land)                                                               \vspace{1.5mm}\\
\textbf{DAC}   & 4 - 7 & \begin{tabular}[c]{@{}l@{}}Energy input, sorbent /\\ solvent materials\end{tabular} & Location-independent         & \begin{tabular}[c]{@{}l@{}}Energy and capital\\ requirements\end{tabular}
\\ \hline
\end{tabular}
\caption{Technological readiness level, cost drivers, benefits and drawbacks of selected NETs}
\end{table}

\subsection{Afforestation and Reforestation (A/R)}
Afforestation and reforestation is the most mature CDR technology due to the extensive human knowledge about forestry, far predating the advent of the issue of climate change through CO$_2$ emissions, and the comparably simple process of restoring forests, either through plantation or natural regeneration. Global forests already sequester 7.6 Gt of CO$_2$ annually without human intervention \parencite{harris_global_2021}. Therefore, A/R benefits from a long history, low technological barriers, and generally shows favorable economics compared to other CDR options \parencite{fuss_negative_2018}. However, A/R is inherently vulnerable to natural disasters such as wildfires, droughts, and pests, which can lead to partial or total reversal of these removals \parencite{wong_can_2003}. This concern is elevated by the increase in such disasters due to climate change \parencite{liang_climate_2025}.\vspace{2mm}

Finally, the sequestration potential of A/R is dependent on environmental conditions. While biomass growth rates and total biomass accumulation per hectare are high in the tropics and sub-tropics, afforestation in temperate regions is significantly less effective. The results in boreal regions are almost an order of magnitude lower than those in the tropics \parencite[4.69]{eggleston_2006_2006}. These differences show the need for spatial differentiation of A/R efforts, and region-focused strategies instead of a blind, global push to maximize reforested areas \parencite{liang_climate_2025}.
\subsection{Enhanced Rock Weathering (ERW)}
Enhanced rock weathering involves spreading finely ground silicate-rich rock material (e.g., from basalt or other olivine-rich rock) on land in order to use and accelerate the natural weathering process of carbon mineralization. It is viewed as a hybrid solution, sitting in between natural and engineered approaches \parencite{manning_enhanced_2025}, and is presently in early-stage laboratory and field tests. Recent field trials, e.g. by the "Carbon Drawdown Initiative", have focused on the technical feasibility of spreading the ground rock and the weathering speed of a variety of different rock and soil combinations and cited mixed conclusions and a need for further research \parencite{carbon_drawdown_initiative_how_2022, hammes_soil_2025}. Similarily, \textcite{madankan_larger_2025} argue that large extraction sites and optimized logistics for processing and transport are essential to make ERW a viable approach.\vspace{2mm}

The primary advantages of ERW over A/R are that it offers highly durable storage in the form of liquid carbonates, and the abundant availability of the input rock materials across the globe \parencite{lackner_progress_1997}. However, \textcite{gaucher_leveraging_2025} argue, while ERW can  improve soil fertility and decrease acidification, the accumulation of heavy metals, their impact on the soil biome, and their potential introduction into the human food chain require careful monitoring. Additionally, and similar to A/R, the weathering process is dependent on the surrounding climate conditions \parencite{hangx_coastal_2009}.

\subsection{Biomass Carbon Removal and Storage (BiCRS)}
Biomass Carbon Removal and Storage is another hybrid approach, and includes different technologies based on the use of biomass. Generally, biomass that sequestered CO$_2$ during its growth is collected and processed into a more permanent form of carbon storage, either through combustion combined with carbon capture and storage (CCS) in the case of Bioenergy with Carbon Capture and Storage (BECCS), or through pyrolysis (thermal decomposition without oxygen) in the cases of Biochar.\vspace{2mm}

BiCRS benefits from a wide range of potential input materials, such as agricultural and forestry waste, animal manure, municipal solid waste, or purpose-grown energy crops \textcite{dipple_building_2021}, and its component technologies are in a relatively mature stage. Additionally, both approaches can deliver significant co-benefits. BECCS generates heat, which can be converted into energy \parencite{hanssen_biomass_2020}, while pyrolysis yields biochar that can be applied to agricultural soil and was found to increase soil microbial biomass, while remaining stable for centuries \parencite{biswal_use_2025}. While BECCS is more suitable for large-scale application, e.g. by refitting coal power plants with CCS and biomass co-firing \parencite{jones_delivering_2021}, Biochar can be deployed in a decentralized fashion using container-sized pyrolysis modules \parencite{horizon_europe_exploiting_2025}.\vspace{2mm}

Compared to ERW, BiCRS generally comes with fewer potential downsides and uncertainties. One risk is the fact that BECCS systems are highly sensitive to emissions along the supply chain and the efficiency of the CCS technology \parencite{jones_delivering_2021}. On the other hand, biochar application to agricultural soil creates some risk of introducing heavy metals and organic toxicants into the food chain, but only at levels not considered health risks \parencite[64]{uba_chancen_2016}. Finally, the use of purpose-grown energy crops always creates a certain amount of biomass competition and land competition with other agricultural applications \parencite[34]{uba_chancen_2016}.
\subsection{Direct Air Capture (DAC)}
Direct Air Capture is a chemical process based on the capture of CO$_2$ directly from the air using solid sorbents (S-DAC) or liquid solvents (L-DAC). Combined with geological storage solutions, it can theoretically provide the highest permanence and easy quantification, allowing for high-quality carbon removal \parencite{stolten_neue_2022}. Additionaly, it is theoretically location-independent, since the only required input is energy. However, in reality, the availability of low-carbon, renewable energy sources may limit suitable geographies for DAC \parencite{eke_comprehensive_2025}.\vspace{2mm}

The main drawback of DAC is the massive energy input necessitated by the thermodynamic limits of seperating CO$_2$ from ambient air, which has relatively low concentrations of carbon dioxide compared to point sources \parencite{lackner_thermodynamics_2013}. Additionally, capital expenses for the construction of DAC plants are also  still high at present, with \textcite{young_cost_2022} suggesting that costs could decrease significantly with increased deployment, similar to the learning effects experienced by solar PV.

\subsection{Other approaches}
There is a wide range of other approaches to carbon removal, such as Ocean Fertilization (OF), Direct Ocean Capture (DOC), and Biomass burial / sinking. However, these methods either suffer from a range of side effects that make verification of impacts extremely hard (such as OF) or are so early in their technological readiness and adoption  that an attempt at modeling them economically would constitute significant guesswork \parencite{smith_state_2024}.

\section{Comparative Economics and Scalability}
\subsection{Costs and Cost Trajectories}
\subsection{Scalability and Resource Constraints}
\section{Social and Political Acceptance}
Social and political acceptance is rarely discussed in CDR literature, which typically focuses on technological and economic feasibility of different NETs. A study by \textcite{karytsas_transnational_2023} found that that social acceptance of CCS (which is not a CDR technology itself, but a part of BECCS) and CO$_2$ pipelines hinges on the perceived safety risks of large scale carbon transport and storage and the trust (or a lack thereof) in their operators. The study also finds that the "Not In My Back Yard" (NIMBY) phenomenon, in the form of resistance to specific projects, is particularly strong for this type of technology.\vspace{2mm}

The socio-political barriers of technologies like ERW are less well understood, but they likely involve concerns over large-scale mining activities and dust polution from rock crushing and spreading \parencite{manning_enhanced_2025, madankan_larger_2025}. Additionally, the political and bureaucratic hurdles of such mining activities are significant. The amount of rock extracted that is required to conduct EWR on a gigaton scale is similar to the size of the entire global coal mining industry, and likely poses the main limiting factor for the deployment of ERW \parencite{hangx_coastal_2009}.\vspace{2mm}

In contrast, the social acceptance of solutions like A/R is much higher, at least as long as it does not compete with agricultural land used for food production \parencite{schirmer_assessing_2014}, potentially owed to the intuitive approach of "just planting trees", but also suggesting some sort of naturalness bias. It may therefore be wise to consider not only the objective net effect of any given NET, but also its subjective perception in the general population. A technology that is perceived as questionable or even dirty may face significant resistance, even if it is objectively beneficial for the global environment at large.